% from ega0-1-fr.tex
\oldpage[0\textsubscript{I}]{78}

pour $M=L=A^m$~; mais alors $\operatorname{Hom}_A(L,N)=N^m$,
$\operatorname{Hom}_A(L,\mathfrak J^{i+1}N)=(\mathfrak J^{i+1}N)^m=
\mathfrak J^{i+1}N^m$ et le résultat est évident. Pour établir la seconde
assertion, notons que l'image de $\operatorname{Hom}_A(M,\mathfrak J^{i+1}N)$ dans
$\operatorname{Hom}_{A_j}(M_j,N_j)$ est nulle pour $j\leq i$, donc $\varphi$ est
continu~; inversement, l'image réciproque dans $\operatorname{Hom}_A(M,N)$ du $0$ de
$\operatorname{Hom}_{A_i}(M_i,N_i)$ est $\operatorname{Hom}_A(M,\mathfrak J^{i+1}N)$, donc $\varphi$
est bicontinu.

Si on suppose seulement que $A$ est un anneau $\mathfrak J$-préadique
noethérien, $M$ et $N$ deux $A$-modules de type fini, séparés pour la
topologie $\mathfrak J$-préadique, la démonstration précédente montre que
la première assertion de (7.8.2) reste valable, et que $\varphi$ est un
isomorphisme topologique de $\operatorname{Hom}_A(M,N)$ sur un sous-module de
$\varprojlim_i\operatorname{Hom}_{A_i}(M_i,N_i)$.

\begin{proposition}[7.8.3]
\label{0.7.8.3-fr}
Sous les hypothèses de (7.8.2), l'ensemble des homomorphismes injectifs
(resp. surjectifs, bijectifs) de $M$ dans $N$ est une partie ouverte de
$\operatorname{Hom}_A(M,N)$.
\end{proposition}

En effet, en vertu de (7.3.5) et (7.1.14), pour que $u$ soit surjectif, il
faut et il suffit que l'homomorphisme correspondant
$u_0:M/\mathfrak JM\to N/\mathfrak JN$ le soit, et l'ensemble des
homomorphismes surjectifs de $M$ dans $N$ est donc image réciproque par
l'application continue $\operatorname{Hom}_A(M,N)\to\operatorname{Hom}_{A_0}(M_0,N_0)$ d'une partie
d'un espace discret. Montrons maintenant que l'ensemble des homomorphismes
injectifs est ouvert~; soit $v$ un tel homomorphisme et posons $M'=v(M)$~;
par le lemme d'Artin--Rees (7.3.2.1), il existe un entier $k\geq 0$ tel que
$M'\cap\mathfrak J^{m+k}N\subset\mathfrak J^mM'$ pour tout $m>0$~; nous
allons voir que pour $w\in\mathfrak J^{k+1}\operatorname{Hom}_A(M,N)$,
$u=v+w$ est injectif, ce qui achèvera la démonstration. En effet, soit
$x\in M$ tel que $u(x)=0$~; prouvons que pour tout $i\geq 0$ la relation
$x\in\mathfrak J^iM$ implique $x\in\mathfrak J^{i+1}M$~; il en résultera
bien
\[
  x\in\bigcap_{i\geq 0}\mathfrak J^iM=(0).
\]
En effet, on a alors $w(x)\in\mathfrak J^{i+k+1}N$, et par ailleurs
$w(x)=-v(x)\in M'$, donc
$v(x)\in M'\cap\mathfrak J^{i+k+1}N\subset\mathfrak J^{i+1}M'$, et comme
$v$ est un isomorphisme de $M$ sur $M'$, $x\in\mathfrak J^{i+1}M$~;
c.q.f.d.

\begin{flushright}
\emph{(A suivre.)}
\end{flushright}
